%方案二

\textbf{动量方法的工程化改进：抑制震荡，加速收敛}

在深度学习训练中，直接使用标准梯度下降法往往会在狭长峡谷型损失曲面中出现剧烈震荡，导致训练速度缓慢甚至难以稳定收敛。为缓解这一问题，动量（Momentum）方法在参数更新中引入“速度”变量，对历史梯度进行指数加权平均，其基本形式为
\[
\mathbf{v}_{k+1} = \beta \mathbf{v}_k + \nabla f(\mathbf{w}_k), \quad
\mathbf{w}_{k+1} = \mathbf{w}_k - \eta \mathbf{v}_{k+1},
\]
其中 $\beta \in (0,1)$ 为动量系数。

从几何直觉上看，动量机制相当于在平缓方向上不断累积“惯性”，从而加速沿峡谷方向的前进；而在陡峭方向上，由于梯度符号频繁交替，动量会相互抵消，从而有效抑制横向震荡。这使得参数更新轨迹由原本的“之”字形震荡，逐渐转变为更加平滑、稳定的下山路径。

在此基础上，RMSProp、Adam 等自适应优化算法进一步结合了梯度平方的滑动平均，实现了对不同方向学习率的自动调节，从工程实践角度显著提升了深度神经网络的训练稳定性与收敛效率